\chapter{Exact four-state entropy (menu 12)}
\label{ch:menu12}
\menulabel{12}

\section{Purpose}

The autoCAS single-orbital entropy itself: for every active orbital, the
four-state quantity $s_i = -\sum_j w_j \ln w_j$ over the occupation weights
(empty, single up, single down, double occupied). Menu~\menu{1} reports a
rigorous \emph{upper bound} for the same number, because an ORCA output alone
cannot supply it (the printed natural occupations are spin-summed). This menu
reconstructs the number from data ORCA does write -- the active-space
Hamiltonian -- and validates the reconstruction against the engine's own
numbers before reporting anything.

\section{What you need}

\begin{itemize}
  \item the \textbf{converged CASSCF output} (the run whose orbitals were
    dumped);
  \item the \textbf{FCIDUMP}: rerun the same job with the converged orbitals
    plus the dump keyword
    (\code{!moread} and \code{\%moinp "previous.gbw"} together with
    \code{!FCIDUMP}). The dump run reports \code{IS NOT FULLY CONVERGED} and
    stops after its first macro-iteration -- that is how the keyword behaves;
    the cross-checks below verify the dumped data against the original run, so
    the message can be ignored for this purpose;
  \item for the \textbf{localized-basis step} (required before reading the
    $0.14$ line): run \code{orca\_loc} on the active window, then
    \code{orca\_2json} on both gbw files (canonical and localized) with a
    configuration requesting the MO coefficients and the $S$ matrix. The two
    exports are the ends of the rotation.
\end{itemize}

\section{Walkthrough}

\lstinputlisting[style=fbkcode, firstline=17, lastline=50,
  title={The menu-12 example session (the N\textsubscript{2} CAS(6,6)
  FCIDUMP chain: converged output, dump, \code{orca\_loc} exports)}]
  {examples/expected/m12_exact_entropy.txt}

The active window is given explicitly here (\code{4 9}); pressing Enter
instead infers it from the occupation table, grows it away from the core when
an active orbital prints as $0.0000$, and validates the result against the
\code{N(occ)=} line.

\section{What you get}

A report (\file{FCIDUMP.fbk.md}, written next to the dump) with: the size of
the determinant space actually solved, the selected root and its $\langle
S^2\rangle$, the cross-checks, both entropy spectra (canonical and localized),
the four weights per localized orbital, the $0.14$-line candidates and the
plateau readout, the preconditions, and the complete citations.

\section{How to read the numbers}

\begin{readbox}
\begin{itemize}
  \item \textbf{The cross-checks come first.} The reconstructed CI energy must
    reproduce the printed CASSCF energy (the gate is $10^{-6}$~Eh; the example
    agrees to $3\times10^{-10}$, limited by the nine decimals of the summary
    line), and the natural occupations must reproduce the \code{N(occ)=} line
    (gate $10^{-4}$). If either fails, the two files do not belong together and
    the menu refuses to continue.
  \item \textbf{Read the threshold in the localized basis only.} The 0.14 line
    and the plateau language are defined for a localized-orbital spectrum. The
    canonical spectrum is printed for orientation; in the example it is small
    throughout ($\le 0.23$), while the localized spectrum shows the real
    structure: four orbitals at $1.355$ (close to $\ln 4$, the maximal mixing
    of the strong $\sigma$ and $\pi$ pairs) and two at $0.37$--$0.39$, with the
    largest relative gap after the fourth. The canonical numbers would have
    suggested ``no multireference character''; the localized numbers say
    otherwise -- the precondition paragraph travels with the report for
    exactly this reason.
  \item \textbf{The four-state density is diagonal for a fixed particle
    number.} That is an exact statement (the $|0\rangle\!\langle 2|$ coherence
    would need a rest configuration with two different electron counts, and
    the spin-flip coherence vanishes for an $S_z$ eigenstate), so the three
    numbers $n_\uparrow$, $n_\downarrow$ and $p_2$ \emph{are} the complete
    one-orbital information -- and they exist only in a reconstruction like
    this. A printed occupation line cannot supply them.
\end{itemize}
\end{readbox}

\section{Worked example}

The session above is the reference chain shipped with the toolkit
(\file{fixtures/orca/n2\_fcidump.*}). To reproduce it for your own system,
run the two steps and the two exports, then feed menu~12; the report's
cross-check lines tell you immediately whether the chain is sound. For an
$f$-block window the determinant count is in the hundreds (the example solves
400), far below the solver's cap of a few thousand per $M_s$ sector.

\section{Caveats, refusals and related sections}

\begin{refusalbox}
\begin{itemize}
  \item A non-converged CASSCF output is refused: the dumped Hamiltonian would
    describe intermediate orbitals and no engine number could validate the
    reconstruction.
  \item The localized step needs \emph{both} exports; one alone is refused
    (the rotation is not defined).
  \item A window whose occupations do not match the printed \code{N(occ)=} is
    refused with a request for the explicit window -- a shifted window of the
    right size cannot pass unnoticed.
  \item The active space is taken as given: size, composition and solution
    branch caveats apply unchanged (menu~\menu{1} checks them). The solver is
    exact within the dumped space but is a dense determinant CI, not a DMRG.
  \item ORCA's own 2-RDM export (\code{orca\_2json} reading \file{.RDM2}
    files) is the documented alternative route, but it was measured
    unreachable for CAS-type methods on ORCA 6.1.1 (six FIC methods, three
    density requests); this FCIDUMP route is the one that works.
\end{itemize}
\end{refusalbox}

Related: \menu{1} (the bound version of this spectrum, and the diagnostics
that should be read alongside), \menu{3} (input generation), and
Chapter~\ref{ch:criteria} for the threshold table.
